%------------------------------
% MICROTYPE
%------------------------------
\microtypecontext{spacing=nonfrench}
%------------------------------
% DIRECTORIES
%------------------------------
\setcounter{secnumdepth}{4}
\setcounter{tocdepth}{4}
%------------------------------
% DATE
%------------------------------
\makeatletter
\newcommand{\monthyeardate}{%
	\iflanguage{english}{\DTMenglishmonthname{\@dtm@month} \@dtm@year}{%
		\DTMgermanmonthname{\@dtm@month} \@dtm@year
	}
}
\makeatother
%------------------------------
% FIGURES, TABLES, LISTS, NUMERATIONS, COUNTER
%------------------------------
\graphicspath{{Figures/}}
\newcounter{includepdfpage}
\newlist{myList}{itemize}{1}
\setlist[myList]{leftmargin=0.0cm,nosep,label={}}
\newcolumntype{s}{>{\collectcell\unit}c<{\endcollectcell}}
\renewcommand{\topfraction}{0.98}
\renewcommand{\bottomfraction}{0.98}
%------------------------------
% DEDICATION
%------------------------------
\newenvironment{dedication_env}
{\Large\begin{quote}\begin{center}\begin{em}}
			{\par\end{em}\end{center}\end{quote}}
%------------------------------
% HEADER
%------------------------------
\lehead{\scshape\leftmark}
\cehead{}
\rehead{}
\lohead{}
\cohead{}
\rohead{\scshape\rightmark}
%------------------------------
% FOOTER
%------------------------------
\lefoot[\pagemark]{\pagemark}
\cefoot{}
\refoot{}
\lofoot{}
\cofoot{}
\rofoot[\pagemark]{\pagemark}
\renewcommand*{\chapterpagestyle}{plain}
%------------------------------
% LISTS OF ...
%------------------------------
\makeatletter
\renewcommand*{\@pnumwidth}{2.5em}
\makeatother
%------------------------------
% LIST OF ABBREVIATIONS
%------------------------------
\newlength{\myabbrevlength}
\setlength{\myabbrevlength}{2cm}
\newlength{\mydescriptionwidth}
\setlength{\mydescriptionwidth}{\linewidth-2\tabcolsep-\myabbrevlength}
\iflanguage{english}{
	\newcommand{\abbrevname}{List of Abbreviations}
}{\iflanguage{ngerman}{
		\newcommand{\abbrevname}{Abkürzungsverzeichnis}}{}
}
\NewAcroTemplate[list]{dbtalongtable}{
	\AcroNeedPackage{array}%
	\def\arraystretch{1.2}
	\acronymsmapF{%
		\AcroAddRow{%
			\leavevmode\acrowrite{short} & \acrowrite{list} \acropagefill \tabularnewline}
	}
	{\AcroRerun}%
	\acroheading%
	\acropreamble%
	\setlength\LTleft{0.5ex}
	\begin{longtable}{>{\bfseries\sffamily}lp{\mydescriptionwidth}}
		\AcronymTable
	\end{longtable}
}
\acsetup{
	format/first-long=\itshape,
	index/use=true,
	make-links=true
}
%------------------------------
% TODOS
%------------------------------
\makeatletter
\iflanguage{ngerman}{
	\renewcommand{\@todonotes@todolistname}{ToDo-Verzeichnis}
}{
	\iflanguage{english}{
		\renewcommand{\@todonotes@todolistname}{List of ToDos}
	}{}
}
\makeatother
%------------------------------
% THEOREM STYLE
%------------------------------
\newtheoremstyle{mytheoremstyle}
{\topsep}
{\topsep}
{\itshape}
{0pt}
{\bfseries}
{}
{5pt plus 1pt minus 1pt}
{\thmname{#1} \thmnumber{#2} \thmnote{\ (#3)}:}
%------------------------------
% COMMANDS FOR THEOREM, ...
%------------------------------
\newcommand{\theoremname}{}
\newcommand{\lemmaname}{}
\newcommand{\examplename}{}
\newcommand{\remarkname}{}
\newcommand{\propositionname}{}
\newcommand{\corollaryname}{}
\iflanguage{english}{
	\renewcommand{\theoremname}{Theorem}
	\renewcommand{\lemmaname}{Lemma}
	\renewcommand{\propositionname}{Proposition}
	\renewcommand{\examplename}{Example}
	\renewcommand{\remarkname}{Remark}
	\renewcommand{\corollaryname}{Corollary}
}{\iflanguage{ngerman}{
		\renewcommand{\theoremname}{Satz}
		\renewcommand{\lemmaname}{Hilfssatz}
		\renewcommand{\examplename}{Beispiel}
		\renewcommand{\propositionname}{Proposition}
		\renewcommand{\remarkname}{Bemerkung}
		\renewcommand{\corollaryname}{Korollar}
	}{}
}
%------------------------------
% SET THEOREM STYLE
%------------------------------
\theoremstyle{mytheoremstyle}
%------------------------------
% THEOREM ENVIRONMENT AND SHARED COUNTER
%------------------------------
\newmdtheoremenv[
linecolor=black,
roundcorner=2pt,
linewidth=2pt
]{theoremenv}{\theoremname}[chapter]

% Unboxed environments share the same chapter-based counter.
% Define them after theoremenv, which owns the counter.
\newtheorem{theorem}[theoremenv]{\theoremname}
\newtheorem{lemma}[theoremenv]{\lemmaname}
\newtheorem{remark}[theoremenv]{\remarkname}
\newtheorem{corollary}[theoremenv]{\corollaryname}
%------------------------------
% LEMMA ENVIRONMENT
%------------------------------
\newmdtheoremenv[
linecolor=blue,
roundcorner=2pt,
linewidth=1pt
]{lemmaenv}[theoremenv]{\lemmaname}
%------------------------------
% EXAMPLE ENVIRONMENT
%------------------------------
\newmdtheoremenv[
linecolor=black,
roundcorner=2pt,
linewidth=0.5pt
]{example}[theoremenv]{\examplename}
%------------------------------
% COROLLARY ENVIRONMENT
%------------------------------
\newmdtheoremenv[
linecolor=black,
roundcorner=2pt,
linewidth=1pt
]{corollaryenv}[theoremenv]{\corollaryname}
%------------------------------
% REMARK ENVIRONMENT
%------------------------------
\newmdtheoremenv[
linecolor=red,
roundcorner=2pt,
linewidth=0.5pt
]{remarkenv}[theoremenv]{\remarkname}
%------------------------------
% DEFINITION ENVIRONMENT
%------------------------------
\newmdtheoremenv[
linecolor=black,
roundcorner=2pt,
linewidth=1pt
]{definition}[theoremenv]{Definition}
%------------------------------
% PROPOSITION ENVIRONMENT
%------------------------------
\newmdtheoremenv[
linecolor=black,
roundcorner=2pt,
linewidth=1pt
]{proposition}[theoremenv]{\propositionname}
%------------------------------
% MATH AND OPERATORS
%------------------------------
\DeclareMathOperator{\rot}{rot}
\newcommand*\diff{\mathop{}\!\textrm{d}}
%------------------------------
% UNITS
%------------------------------
\newcommand{\SIRangePhrase}{}
\sisetup{
	detect-family=true,
	detect-mode=true,
	mode=text,
	range-units=single,
	retain-explicit-plus,
	retain-unity-mantissa=false,
	retain-zero-exponent=false,
	separate-uncertainty=true,
}
\DeclareSIUnit\torr{Torr}
\DeclareSIUnit\fahrenheit{\degree F}
\DeclareSIUnit\atm{atm}
\DeclareSIUnit\anno{a}
\DeclareSIUnit\bar{bar}
\DeclareSIUnit\kJmol{\kilo\joule\per\mole}
\DeclareSIUnit\kJmolK{\kilo\joule\per\mole\per\kelvin}
%------------------------------
% CHEMISTRY
%------------------------------
\chemsetup{%
	formula=chemformula,
	greek=upgreek,
}
%------------------------------
% COLORS
%------------------------------
\definecolor{dbta_blue}{RGB}{51,102,153}
\definecolor{tub_red}{RGB}{197,14,31}
%------------------------------
% XPARSE
%------------------------------
\NewDocumentCommand{\pder}{ O{} O{} m }{\dfrac{\partial^{#2}#1}{\partial#3^{#2}}}
\NewDocumentCommand{\roundbrack}{ m }{\left(#1 \right)}
\NewDocumentCommand{\squarebrack}{ m }{\left[#1 \right]}
\NewDocumentCommand{\ind}{m O{}}{\mathrm{#1}\mathnormal{,#2}}
%------------------------------
% CODE
%------------------------------
\renewcommand{\lstlistingname}{Code}
%------------------------------
% ENUMERATE ENVIRONMENT
%------------------------------
\renewcommand{\labelenumi}{\arabic{enumi})}
\renewcommand{\labelenumii}{\arabic{enumi}.\arabic{enumii}}
\renewcommand{\labelenumiii}{\arabic{enumi}.\arabic{enumii}.\arabic{enumiii}}
\setlist[enumerate,1]{%
	before*=\onehalfspacing,
	itemsep=0.15\baselineskip,
	labelindent=\parindent,
}%
%------------------------------
% ITEMIZE ENVIRONMENT
%------------------------------
\setlist[itemize,1]{%
	label={--},
	before*=\onehalfspacing,
	itemsep=0.15\baselineskip,
	labelindent=\parindent,
}%
%------------------------------
% GERMAN CAPTION REDEFINITIONS
%------------------------------
\addto\captionsngerman{%
	\renewcommand{\figurename}{Abb.}
	\renewcommand{\tablename}{Tab.}
	\SetAlgorithmName{Algorithmus}{Algorithmus}{Algorithmenverzeichnis}
	\renewcommand{\nomname}{Symbolverzeichnis}
	\renewcommand{\lstlistlistingname}{Codeverzeichnis}
	\def\indexname{Indexverzeichnis}
}
%------------------------------
% ENGLISH CAPTION REDEFINITIONS
%------------------------------
\addto\captionsenglish{
	\renewcommand{\figurename}{Fig.}
	\renewcommand{\tablename}{Tab.}
	\renewcommand{\nomname}{List of Symbols}
	\renewcommand{\lstlistlistingname}{List of Codes}
	\def\indexname{Index}
}
%------------------------------
% ADD TOCLIST ALGORITHM
%------------------------------
\addtotoclist[algorithm2e]{loa}
\setuptoc{loa}{totoc}
\newcommand*{\listofloaname}{\listalgorithmcfname}
\renewcommand{\listofalgorithms}{\listoftoc{loa}}
%------------------------------
% AUTOREF GERMAN
%------------------------------
\addto\extrasngerman{%
	\def\chapterautorefname{Kap.}%
	\def\sectionautorefname{Abschn.}%
	\def\subsectionautorefname{Abschn.}%
	\def\subsubsectionautorefname{Abschn.}%
	\def\figureautorefname{Abb.}%
	\def\tableautorefname{Tab.}%
	\def\theoremenvautorefname{Satz}%
	\def\lemmaenvautorefname{Hilfssatz}%
	\def\exampleautorefname{Bsp.}%
	\def\exampleautorefnamelong{Beispiel}%
	\def\corollaryenvautorefname{Korollar}%
	\def\remarkenvautorefname{Bem.}%
	\def\definitionautorefname{Def.}%
	\def\algorithmautorefname{Alg.}%
	\def\assumptionautorefname{Ann.}%
	\def\equationautorefname~#1\null{Gl.~(#1)\null}%
	\def\chapterautorefnamelong{Kapitel}%
	\def\sectionautorefnamelong{Abschnitt}%
	\def\subsectionautorefnamelong{Abschnitt}%
	\def\subsubsectionautorefnamelong{Abschnitt}%
	\def\figureautorefnamelong{Abbildung}%
	\def\tableautorefnamelong{Tabelle}%
	\def\equationautorefnamelong~#1\null{Gleichung~(#1)\null}%
	\def\theoremenvautorefnamelong{Satz}%
	\def\lemmaenvautorefnamelong{Hilfssatz}%
	\def\corollaryenvautorefnamelong{Korollar}%
	\def\remarkenvautorefnamelong{Bemerkung}%
	\def\definitionautorefnamelong{Definition}%
	\def\algorithmautorefnamelong{Algorithmus}%
	\renewcommand{\SIRangePhrase}{~bis~}%
	\sisetup{%
		locale = DE,%
		range-phrase={\SIRangePhrase},
	}%
}
%------------------------------
% AUTOREF ENGLISH
%------------------------------
\addto\extrasenglish{%
	\def\chapterautorefname{Chap.}%
	\def\sectionautorefname{Sec.}%
	\def\subsectionautorefname{Sec.}%
	\def\subsubsectionautorefname{Sec.}%
	\def\figureautorefname{Fig.}%
	\def\tableautorefname{Tab.}%
	\def\exampleautorefname{Example}%
	\def\exampleautorefnamelong{Example}%
	\def\equationautorefname~#1\null{Eq.~(#1)\null}%
	\def\theoremenvautorefname{Theorem}%
	\def\lemmaenvautorefname{Lemma}%
	\def\corollaryenvautorefname{Corollary}%
	\def\remarkenvautorefname{Remark}%
	\def\definitionautorefname{Def.}%
	\def\algorithmautorefname{Alg.}%
	\def\assumptionautorefname{Ass.}%
	\def\chapterautorefnamelong{Chapter}%
	\def\sectionautorefnamelong{Section}%
	\def\subsectionautorefnamelong{Section}%
	\def\subsubsectionautorefnamelong{Section}%
	\def\figureautorefnamelong{Figure}%
	\def\tableautorefnamelong{Table}%
	\def\equationautorefnamelong~#1\null{Equation~(#1)\null}%
	\def\theoremenvautorefnamelong{Theorem}%
	\def\lemmaenvautorefnamelong{Lemma}%
	\def\corollaryenvautorefnamelong{Corollary}%
	\def\remarkenvautorefnamelong{Remark}%
	\def\definitionautorefnamelong{Definition}%
	\def\algorithmautorefnamelong{Algorithm}%
	\renewcommand{\SIRangePhrase}{~to~}%
	\sisetup{%
		locale = US,%
		range-phrase={\SIRangePhrase},
	}%
}
%------------------------------
% AUTOREF (LONG FORM)
%------------------------------
\newcommand{\Autoref}[1]{%
	\begingroup%
	\renewcommand{\chapterautorefname}{\chapterautorefnamelong}%
	\renewcommand{\sectionautorefname}{\subsectionautorefnamelong}%
	\renewcommand{\subsectionautorefname}{\subsectionautorefnamelong}%
	\renewcommand{\subsubsectionautorefname}{\subsectionautorefnamelong}%
	\renewcommand{\exampleautorefname}{\exampleautorefnamelong}%
	\renewcommand{\figureautorefname}{\figureautorefnamelong}%
	\renewcommand{\tableautorefname}{\tableautorefnamelong}%
	\renewcommand{\equationautorefname}{\equationautorefnamelong}%
	\renewcommand{\theoremenvautorefname}{\theoremenvautorefnamelong}%
	\renewcommand{\corollaryenvautorefname}{\corollaryenvautorefnamelong}%
	\renewcommand{\lemmaenvautorefname}{\lemmaenvautorefnamelong}%
	\renewcommand{\remarkenvautorefname}{\remarkenvautorefnamelong}%
	\renewcommand{\definitionautorefname}{\definitionautorefnamelong}%
	\renewcommand{\algorithmautorefname}{\algorithmautorefnamelong}%
	\renewcommand{\assumptionautorefname}{\assumptionautorefnamelong}%
	\autoref{#1}%
	\endgroup%
}
%------------------------------
% REFERENCES AND CITATIONS
%------------------------------
\renewcaptionname{english}{\bibname}{List of References}
\renewcaptionname{ngerman}{\bibname}{Literaturverzeichnis}
\setlength\bibitemsep{0.5\baselineskip}
\setlength{\bibhang}{5ex}
\DeclareFieldFormat[
article,
inbook,
incollection,
inproceedings,
patent,
thesis,
unpublished
]{citetitle}{#1\midsentence\isdot}
\DeclareFieldFormat[
article,
inbook,
incollection,
inproceedings,
patent,
thesis,
unpublished
]{title}{#1\midsentence\isdot}
\DeclareFieldFormat[
article,
]{titlecase:title}{\MakeSentenceCase*{#1}}
\DeclareFieldFormat[
inproceedings,
]{booktitle}{\itshape #1}
\renewcommand*{\nameyeardelim}{\addcomma\space}
\DeclareFieldFormat{pages}{#1}
\DeclareFieldFormat{type}{\itshape #1}
\DefineBibliographyStrings{english}{
	backrefpage = {cit. on p.},
	backrefpages = {cit. on pp.},
}
\DefineBibliographyStrings{ngerman}{
	backrefpage = {zit. auf S.},
	backrefpages = {zit. auf den S.},
}
\renewbibmacro{in:}{%
	\ifboolexpr{%
		test {\ifentrytype{article}}%
		or
		test {\ifentrytype{book}}%
		or
		test {\ifentrytype{inproceedings}}%
	}{}{\printtext{\bibstring{in}\intitlepunct}}%
}
\DefineBibliographyStrings{ngerman}{
	urlseen = {letzter Zugriff},
}
\DefineBibliographyStrings{english}{
	urlseen = {last access},
}
\renewbibmacro*{volume+number+eid}{
	\printfield{volume}
	\printfield[parens]{number}%
	\setunit{\addcomma\space}%
	\printfield{eid}
}
\DeclareDelimFormat[bib]{nametitledelim}{\addcolon\space}
\DeclareDelimFormat{finalnamedelim}{\addcomma\space}
\DeclareDelimFormat[textcite]{finalnamedelim}{
	\addspace\&\space}
\DeclareDelimFormat[parencite]{finalnamedelim}{
	\addspace\bibstring{and}\space}

\renewbibmacro*{begentry}{%
	\iffieldundef{shorthand}{}{%
		\printfield{shorthand}\addspace ---\addspace
	}
}
\DeclareSortingTemplate{nyt}{
	\sort{%
		\field{presort}%
	}%
	\sort[final]{%
		\field{sortkey}%
	}%
	\sort{%
		\field{shorthand}%
		\field{author}%
		\field{editor}%
		\field{title}%
	}%
	\sort{%
		\field{year}
	}%
	\sort{%
		\field{title}%
	}%
}

\renewbibmacro*{textcite}{%
	\ifnameundef{labelname}{%
		\iffieldundef{shorthand}{%
			\usebibmacro{cite:label}%
			\setunit{%
				\global\booltrue{cbx:parens}%
				\printdelim{nameyeardelim}\bibopenparen%
			}%
			\ifnumequal{%
				\value{citecount}%
			}{1}{%
				\usebibmacro{prenote}%
			}{}%
			\usebibmacro{cite:labeldate+extradate}%
		}{%
			\usebibmacro{cite:shorthand}%
	}}{%
		\iffieldundef{shorthand}{%
			\printnames{labelname}%
			\setunit{%
				\global\booltrue{cbx:parens}%
				\printdelim{nameyeardelim}\bibopenparen%
			}%
			\ifnumequal{%
				\value{citecount}%
			}{1}{%
				\usebibmacro{prenote}%
			}{}%
			\usebibmacro{citeyear}%
		}{%
			\ifciteseen{%
				\usebibmacro{cite:shorthand}%
			}{%
				\printnames{labelname}%
				\setunit{%
					\global\booltrue{cbx:parens}%
					\printdelim{nameyeardelim}\bibopenparen%
				}%
				\ifnumequal{%
					\value{citecount}%
				}{1}{%
					\usebibmacro{prenote}%
				}{}%
				\iffieldundef{labelyear}{%
					\usebibmacro{cite:label}%
					\setunit{%
						\addcomma%
						\addspace%
					}%
					\iffieldundef{shorthandintro}{%
						\printtext{\bibstring{citedas}}%
					}{%
						\printtext{\printfield{shorthandintro}}%
					}%
					\addspace%
					\usebibmacro{cite:shorthand}%
				}{%
					\usebibmacro{cite:labeldate+extradate}%
					\setunit{%
						\addcomma%
						\addspace%
					}%
					\iffieldundef{shorthandintro}{%
						\printtext{\bibstring{citedas}}%
					}{%
						\printtext{\printfield{shorthandintro}}%
					}%
					\addspace%
					\usebibmacro{cite:shorthand}}}}}}%
%------------------------------
% ONLY FOR DISSERTATIONS
%------------------------------
\DeclareBibliographyCategory{inbib}
\newcommand{\printpublication}[1]{\AtNextCite{\defcounter{maxnames}{99}}\fullcite{#1}}
\makeatletter
\AtEveryCitekey{%
	\ifcsstring{blx@delimcontext}{fullcite}
	{}
	{\addtocategory{inbib}{\thefield{entrykey}}}}
\makeatother
%------------------------------
% INTERNET
%------------------------------
\renewcommand*{\UrlFont}{\upshape}
%------------------------------
% HYPERREF
%------------------------------
\hypersetup{%
	bookmarksnumbered=true,
	bookmarksopen=false,
	citebordercolor=blue,
	colorlinks=false,
	hypertexnames=true,
	linkbordercolor=red,
	linktocpage,
	pdfpagelayout=TwoPageRight,
	urlbordercolor=cyan,
}%
%------------------------------
% MAKEINDEX
%------------------------------
\makenomenclature
\iflanguage{ngerman}{
	\newcommand{\indextitle}{Indexverzeichnis}
}{
	\iflanguage{english}{
		\newcommand{\indextitle}{Index}
	}{}
}
\makeindex[
columns=2,
intoc=true,
options= -s 05_Literature_and_Index/myindexstyle.ist,
title={\indextitle},
]
%------------------------------
% OWN COMMANDS
%------------------------------
% ---------- commonly used sets ----------
\DeclareMathOperator{\N}{\mathbb{N}}
\DeclareMathOperator{\C}{\mathbb{C}}

% ---------- Category theory operators ----------
\DeclareRobustCommand{\CatC}{\ensuremath{\mathcal{C}}}
\DeclareRobustCommand{\CatD}{\ensuremath{\mathcal{D}}}
\DeclareRobustCommand{\Cat}[1]{\ensuremath{\mathcal{#1}}}
\DeclareMathOperator{\Ob}{Ob}
\NewDocumentCommand{\adj}{mm}{\ensuremath{#1 \dashv #2}}

% ---------- (Pre)Sheaves ----------
\DeclareRobustCommand{\Psh}[1]{\ensuremath{\cat{PSh}\!\left(#1\right)}}
\DeclareRobustCommand{\PshC}{\Psh{\CatC}}
\DeclareRobustCommand{\Sh}[1]{\ensuremath{\cat{Sh}\!\left(#1\right)}}
\DeclareRobustCommand{\ShC}{\Sh{\CatC}}
\DeclareRobustCommand{\PshSet}[1]{\ensuremath{\Set^{\Opp{#1}}}}
\DeclareRobustCommand{\PshSetC}{\PshSet{\CatC}}

% ---------- Opposite category ----------
\DeclareRobustCommand{\op}{\mathrm{op}}
\NewDocumentCommand{\Opp}{m}{\ensuremath{#1^{\op}}}
\DeclareRobustCommand{\CatCop}{\Opp{\CatC}}
\DeclareRobustCommand{\CatDop}{\Opp{\CatD}}

% ---------- Categories (sans serif) ----------
\DeclareRobustCommand{\cat}[1]{\ensuremath{\mathsf{#1}}}
\DeclareRobustCommand{\Cat}{\cat{Cat}}
\DeclareRobustCommand{\Sch}{\cat{Sch}}
\DeclareRobustCommand{\Set}{\cat{Set}}
\DeclareRobustCommand{\Top}{\cat{Top}}
\DeclareRobustCommand{\Grp}{\cat{Grp}}
\DeclareRobustCommand{\Grpd}{\cat{Grpd}}
\DeclareRobustCommand{\Alg}{\cat{Alg}}
\DeclareRobustCommand{\Ab}{\cat{Ab}}
\DeclareRobustCommand{\Ring}{\cat{Ring}}
\DeclareRobustCommand{\CRing}{\cat{CRing}}
\DeclareRobustCommand{\Mod}{\cat{Mod}}

% ---------- Limits and colimits ----------
\DeclareMathOperator*{\colim}{colim}
\DeclareMathOperator*{\Lim}{lim}
\DeclareMathOperator*{\Eq}{Eq}
\DeclareMathOperator*{\Coeq}{Coeq}

% ---------- Over/under categories ----------
\NewDocumentCommand{\Over}{mm}{\ensuremath{#1_{/#2}}}
\NewDocumentCommand{\Under}{mm}{\ensuremath{#1_{#2/}}}
\NewDocumentCommand{\SchOver}{m}{\Over{\Sch}{#1}}
\NewDocumentCommand{\SchUnder}{m}{\Under{\Sch}{#1}}

% ---------- Algebraic geometry operators ----------
\DeclareMathOperator{\Spec}{Spec}
\DeclareMathOperator{\Proj}{Proj}
\DeclareMathOperator{\Hom}{Hom}
\DeclareMathOperator{\End}{End}
\DeclareMathOperator{\Aut}{Aut}
\NewDocumentCommand{\Specof}{m}{\ensuremath{\Spec\!\left(#1\right)}}
\NewDocumentCommand{\Homof}{mm}{\ensuremath{\Hom\!\left(#1,#2\right)}}

% ---------- Structure sheaf ----------
\DeclareRobustCommand{\OO}{\mathscr{O}}
\DeclareRobustCommand{\OOX}{\mathscr{O}_X}
\NewDocumentCommand{\res}{mm}{#1\!\mid_{#2}}

% ---------- Algebraic spaces ----------
\DeclareRobustCommand{\AlgSp}{\cat{AlgSp}}
\NewDocumentCommand{\AlgSpOver}{m}{\Over{\AlgSp}{#1}}
\NewDocumentCommand{\AlgSpUnder}{m}{\Under{\AlgSp}{#1}}

% ---------- Quasi-coherent sheaves ----------
\DeclareRobustCommand{\QCoh}{\cat{QCoh}}
\NewDocumentCommand{\QCohof}{m}{\ensuremath{\QCoh\!\left(#1\right)}}
\DeclareRobustCommand{\QCohAlg}{\cat{QCohAlg}}
\NewDocumentCommand{\QCohAlgof}{m}{\ensuremath{\QCohAlg\!\left(#1\right)}}
\DeclareRobustCommand{\QCohX}{\QCohof{X}}
\DeclareRobustCommand{\QCohY}{\QCohof{Y}}

% ---------- Stacks ----------
\DeclareRobustCommand{\Stk}{\cat{Stk}}
\NewDocumentCommand{\StkOver}{m}{\Over{\Stk}{#1}}
\NewDocumentCommand{\StkUnder}{m}{\Under{\Stk}{#1}}
\NewDocumentCommand{\Stacksof}{m}{\ensuremath{\Stk\!\left(#1\right)}}
\DeclareRobustCommand{\AlgStk}{\cat{AlgStk}}
\NewDocumentCommand{\AlgStkOver}{m}{\Over{\AlgStk}{#1}}
\NewDocumentCommand{\AlgStkUnder}{m}{\Under{\AlgStk}{#1}}

% ---------- Infinity-stacks ----------
\DeclareRobustCommand{\inftyGrpd}{\ensuremath{\cat{\infty\text{-}Grpd}}}
\DeclareRobustCommand{\inftyStk}{\ensuremath{\cat{\infty\text{-}Stk}}}
\NewDocumentCommand{\inftyStacksof}{m}{\ensuremath{\inftyStk\!\left(#1\right)}}
\NewDocumentCommand{\inftyStkOver}{m}{\Over{\inftyStk}{#1}}
\NewDocumentCommand{\inftyStkUnder}{m}{\Under{\inftyStk}{#1}}

% ---------- Affine schemes ----------
\DeclareRobustCommand{\Aff}{\cat{Aff}}
\NewDocumentCommand{\AffOver}{m}{\Over{\Aff}{#1}}
\NewDocumentCommand{\AffUnder}{m}{\Under{\Aff}{#1}}

% ---------- Homological algebra ----------
\DeclareMathOperator{\Tot}{Tot}
\NewDocumentCommand{\Totof}{m}{\ensuremath{\Tot\!\left(#1\right)}}
\DeclareMathOperator{\Desc}{Desc}

% ---------- Frequently used terms ----------
\DeclareRobustCommand{\etale}{\'etal\'e}
\DeclareRobustCommand{\Etale}{\'Etal\'e}
